\HLASourcePage{31}{30}
\thispagestyle{hla-body}
\begin{hlafixedpage}
  \HLADefinitionAt{69.79}{%
    \HLADefinitionAnchor{pair}\HLADefinitionTerm{pair}: A grouping of two related elements (a first component and a second component), the combination of\linebreak
    which is treated as an entity. An example of a pair would be an instance attribute grouped with its current\linebreak
    value.%
  }

  \HLADefinitionAt{117.78}{%
    \HLADefinitionAnchor{parameter}\HLADefinitionTerm{parameter}: A named characteristic of an interaction.%
  }

  \HLADefinitionAt{141.78}{%
    \HLADefinitionAnchor{passel}\HLADefinitionTerm{passel}: A group of attribute handle/value pairs from an \emph{Update Attribute Values} service invocation that are\linebreak
    delivered together via a \emph{Reflect Attribute Values \textdagger{}} service invocation. All pairs within the passel have the\linebreak
    same user-supplied tag, sent message order type, transportation type, receive message order type, timestamp\linebreak
    (if present), and set of sent region designators (if present). A passel is a message.%
  }

  \HLADefinitionAt{201.78}{%
    \HLADefinitionAnchor{passive-subscription}\HLADefinitionTerm{passive subscription}: A request to the runtime infrastructure (RTI) for the kinds of data (object classes and\linebreak
    attributes as well as interactions) that the joined federate is currently interested in receiving. However, unlike\linebreak
    an active subscription, this information is not used by the RTI to arrange for data to be delivered, nor is it\linebreak
    used to tell publishing joined federates that another joined federate is subscribing to that data (by way of\linebreak
    \emph{Start/Stop Registration}, \emph{Turn On/Off Updates}, or \emph{Turn On/Off Interactions} service invocations). This form of\linebreak
    subscription is provided to support certain types of logger joined federates.%
  }

  \HLADefinitionAt{285.78}{%
    \HLADefinitionAnchor{promoted}\HLADefinitionTerm{promoted}: Pertaining to an object instance, as known by a particular joined federate: having a discovered\linebreak
    class that is a superclass of its registered class.%
  }
  \HLADefinitionNoteAt{90}{321.98}{NOTE---See 5.1.3.}

  \HLADefinitionAt{343.81}{%
    \HLADefinitionAnchor{publish}\HLADefinitionTerm{publish}: To announce to a federation the information a federate may provide to the federation.%
  }

  \HLADefinitionAt{367.81}{%
    \HLADefinitionAnchor{published}\HLADefinitionTerm{published}:%
  }

  \HLAFixedAlphaList{100.02}{391.80}{421.98}{8}{%
    \item Pertaining to an object class from the perspective of a given joined federate: having at least one\linebreak
      available attribute of the object class that was an argument, along with the object class, to a \emph{Publish}\linebreak
      \emph{Object Class Attributes} service invocation that was not subsequently unpublished via the \emph{Unpublish}\linebreak
      \emph{Object Class Attributes} service.\par\vspace{9.4pt}{\fontsize{9pt}{12pt}\selectfont NOTE---See 5.1.2.}
    \item Pertaining to an interaction class from the perspective of a given joined federate: being an argument\linebreak
      to a \emph{Publish Interaction Class} service invocation that was not subsequently followed by an\linebreak
      \emph{Unpublish Interaction Class} service invocation for that interaction class.\par\vspace{9.4pt}{\fontsize{9pt}{12pt}\selectfont NOTE---See 5.1.3.}
  }

  \HLADefinitionAt{537.79}{%
    \HLADefinitionAnchor{published-attributes-of-an-object-class}\HLADefinitionTerm{published attributes of an object class}: The class attributes that have been arguments to \emph{Publish Object}\linebreak
    \emph{Class Attributes} service invocations by a given joined federate for that object class that have not\linebreak
    subsequently been unpublished (either individually or by unpublishing the whole object class), and possibly\linebreak
    the \textbf{HLAprivilegeToDeleteObject} attribute for that object class.%
  }
  \HLADefinitionNoteAt{90}{597.98}{NOTE---See 5.1.2.}

  \HLADefinitionAt{619.81}{%
    \HLADefinitionAnchor{range}\HLADefinitionTerm{range}: A subset of a dimension, defined by an ordered pair of values, the first being the range lower bound\linebreak
    and the second being the range upper bound. This pair of values defines a semi-open interval [range lower\linebreak
    bound, range upper bound) (i.e., the range lower bound is the smallest member of the interval, and the range\linebreak
    upper bound is just greater than any member of the interval).%
  }
  \HLADefinitionNoteAt{90}{680.00}{NOTE---See 9.1.1.}
\end{hlafixedpage}
\clearpage
